\documentclass[11pt,a4paper]{article}
\usepackage[T1]{fontenc}
\usepackage[utf8]{inputenc}
\usepackage[portuguese]{babel}
\usepackage[margin=2.3cm]{geometry}
\usepackage{enumitem}
\usepackage{titlesec}
\usepackage{xcolor}
\usepackage[hidelinks]{hyperref}
\usepackage{booktabs}
\usepackage{amssymb}

\definecolor{cinza}{gray}{0.35}
\titleformat{\section}{\large\bfseries}{}{0pt}{}[\vspace{-6pt}\rule{\textwidth}{0.4pt}]
\titlespacing{\section}{0pt}{14pt}{6pt}
\newcommand{\lin}[2]{\noindent\textbf{#1}\hfill{\color{cinza}\small #2}\par}
\setlist[itemize]{leftmargin=1.2em,itemsep=1pt,topsep=2pt}
\pagestyle{empty}

\begin{document}

% ─── Licença ──────────────────────────────────────────────────────────────
% CC BY 4.0 · Copyright (c) 2026 Aarão Melo Lopes · ver LICENSE na raiz.


\begin{center}
{\LARGE\bfseries Aarão Melo Lopes}\\[3pt]
{\color{cinza}Engenheiro eletricista · matemática aplicada e sistemas}\\[2pt]
{\color{cinza}\small e um mundo narrativo construído --- \emph{procuro parceiros dos dois lados}}\\[2pt]
{\small aarao.melo.lopes@gmail.com}
\end{center}

\vspace{-4pt}
\section*{O que faço}

Construo uma teoria algébrica e as ferramentas que a medem. O trabalho tem uma pergunta só ---
\emph{o que é uma unidade} --- e dela derivam-se as operações, a ordem, o tempo e a geometria. Cada
afirmação vem acompanhada de um medidor em C que a pode desmentir.

\medskip
\noindent Isto não é um projeto recente. Começou num texto de números hipercomplexos escrito na
graduação, atravessou o mestrado no IME-USP, e é hoje um corpo de trabalho com três documentos e
seis repositórios. \textbf{Procuro parceiros para o continuar} --- em matemática, em engenharia,
ou em quem queira usá-lo.

\medskip
\noindent \textbf{E há um teste que este trabalho passou sem o ter procurado.} Há meses que
tenho na mão um ficheiro de $33$ linhas --- um contador de bits escrito em julho, antes de a
teoria ter forma --- e que não sabia fundamentar. \emph{A derivação chegou a ele sozinha}: o par de
máscaras complementares, a involução que as troca, a leitura sem estado, e uma saída que devolve
\textbf{a soma e uma das partes} em vez das duas --- que é exatamente a coordenada que mede e a que
ordena. Nada disso foi posto lá pela teoria; foi reencontrado por ela, e está medido propriedade a
propriedade. \emph{O que se fez depois não foi construir aquilo: foi provar que era aquilo.}

\medskip
\noindent \emph{E não é só papel.} A teoria produziu uma máquina com \textbf{uma instrução e
nenhum estado} --- sem \emph{runtime}, o binário não tem sequer a secção \texttt{.bss}. E a mesma
estrutura já existia em silício: a tabela de comutação de um inversor de três níveis \emph{é} o
corpo mínimo da teoria. \textbf{A proposta I} descreve-o.

\section*{Produção}

\lin{Uma teoria da dualidade, em três volumes}{2026, em curso}
\begin{itemize}
\item \textbf{Teoria} (58 pp.) --- uma lei única e as suas interpretações. \textbf{21 enunciados
e 21 provas}, um por um: 18 teoremas, 2 proposições e 1 corolário, \textbf{nenhum sem
demonstração}. Uma única definição em todo o volume --- a do corpo dual --- e tudo o resto
derivado dela.
\item \textbf{Catálogo} (446 pp.) --- o bestiário dos corpos algébricos, cada um com a sua régua e
a sua dinâmica, e \textbf{275 medidores} citados nominalmente, em 357 chamadas.
\item \textbf{Enredo} (360 pp.) --- a mesma estrutura contada como narrativa, sem uma fórmula.
\item \textbf{306 medidores em C} --- a bateria corre inteira e \textbf{sai a zero}: 306 verdes,
0 falhas. \textbf{115 deles não usam um único número de vírgula flutuante}; a aritmética é inteira e
o resultado é exacto ou não é. A bateria recusa-se a passar quando uma afirmação é falsa, e
\emph{vários destes medidores já desmentiram o autor.}
\end{itemize}

\medskip
\lin{\emph{O sistema dos números hipercomplexos} --- livro próprio}{graduação}
Definição de $\mathbb{H}$ sobre $\mathbb{R}^n$ com produto próprio, associador, conjugação, norma
multiplicativa e forma polar com De Moivre por indução. \textbf{É a semente do trabalho atual}: a
identidade $\hat z^{\,2}=-1$ e a inversa $z^{-1}=\bar z/|z|^2$ escritas aí são, hoje, as duas leis
da teoria.

\medskip
\lin{\emph{Álgebra Bilinear} --- notas próprias}{graduação}
Formas e produtos bilineares, base ordenada, e a representação $f(u,v)=[u]^{T}_{B}A[v]_{B}$ --- a
matriz como leitura de uma operação numa base, que é o teorema central do trabalho atual.

\medskip
\lin{Seis sistemas construídos, 3\,954 registos de versão}{abr.\ 2026 -- ago.\ 2026}
Cada um começou antes de o anterior terminar; sobrepõem-se todos --- e todos existem em disco,
com a contagem de commits ao lado: \texttt{hiper} (1\,117), \texttt{chess} (819), \texttt{tiffany}
(902), \texttt{broca-so} (688), \texttt{machine} (359), \texttt{kernel} (69). Compilador
multi-formato, máquina virtual de pilha, núcleo mínimo, sistema operativo com orquestração de
agentes, motor de xadrez, e o corpo teórico atual.
Publicação automatizada com portão de verificação --- nada sobe sem os documentos compilarem e a
bateria passar.

\section*{Formação}

\lin{Mestrado em Ciência da Computação --- IME-USP}{2020 -- 2022}
Orientador: Prof.\ Dr.\ Ronaldo Fumio Hashimoto. \textbf{Bolsista CAPES / Demanda Social.}
\textbf{72 créditos} em disciplinas --- o programa exige 44 --- com \textbf{conceito máximo em oito
das nove}. Dissertação não depositada; a nona disciplina, e o que ela abriu, estão na
\emph{Proposta de trabalho} adiante.

\begin{itemize}
\item MAC5711 Análise de Algoritmos \textbf{·} MAC5770 Introdução à Teoria dos Grafos
\item MAC6995 Representações Hierárquicas (morfologia matemática) \textbf{·} MAC6903 Processamento
de Imagens usando Grafos \textbf{·} MAC5768 Visão e Processamento de Imagens
\item MAC6916 Modelos Probabilísticos Baseados em Grafos \textbf{·} MAC5742 Computação Paralela e
Distribuída \textbf{·} MAC5744 Introdução à Computação Gráfica
\end{itemize}

\medskip
\lin{Extensão universitária --- IME-USP}{verão de 2021}
\textbf{Cálculo no $\mathbb{R}^n$} e \textbf{Espaços Métricos}, 60 h cada, aprovado com
\textbf{100\,\% de frequência} em ambos. Cursados por iniciativa própria, fora do currículo
obrigatório.

\medskip
\lin{Bacharelado em Engenharia Elétrica --- UFRR}{2011 -- 2018}
Média de conclusão \textbf{8{,}66} · 4\,201 horas integralizadas.

\smallskip
\noindent\textbf{Nota máxima (10{,}0) em:} Cálculo I \textbf{·} Cálculo Diferencial e Integral I
\textbf{·} Eletromagnetismo I \textbf{·} Física Experimental III \textbf{·} Conversão de Energia II
\textbf{·} Fundamentos de Economia. \\
\textbf{Acima de 9:} Estatística (9{,}73) \textbf{·} Química Geral (9{,}89) \textbf{·} Cálculo
Numérico (9{,}5) \textbf{·} Métodos Matemáticos (9{,}5) \textbf{·} Geometria Analítica (9{,}49)
\textbf{·} Álgebra Linear (9{,}17) \textbf{·} Geração de Energia (9{,}9).

\section*{Trabalho de conclusão e docência}

\lin{\emph{Controle Direto de Torque do Motor de Indução Trifásico}}{UFRR, 2018}
123 pp. Orientadora: Prof.\textsuperscript{a} Dr.\textsuperscript{a} Susset Guerra Jiménez.
\textbf{Nota 10 dos três membros da banca}, com recomendação escrita em ata: \emph{«a banca sugere
publicação de artigos».} Modulação por vetores espaciais, inversores multinível, tração veicular.

\smallskip
\noindent \textbf{E oito anos depois a teoria reencontrou-o, por outro caminho.} A tabela de
comutação que projetei nesse trabalho \emph{é} o corpo mínimo da teoria: cada fase toma
$\{-1,0,+1\}$, as três dão $3^3=27$ estados, $19$ vetores e \textbf{seis pares redundantes} --- e
a redundância, que ali servia para equilibrar o ponto médio, \emph{é} a dualidade. O DTC decide
pela fase e pela amplitude do erro, \textbf{sem modulador e sem histórico}: é a máquina sem memória
da \textbf{Proposta I}, e eu tinha-a construído antes de ter o formalismo que a explica.
\emph{Derivei-a por teorema em 2026; tinha-a projetado em 2018 porque o modulador custava tempo.}

\medskip
\lin{Monitoria --- 473 horas}{UFRR, 2013 e 2016--17}
\textbf{Cálculo Diferencial e Integral I} (Engenharia Elétrica, 108 h) e \textbf{Álgebra Linear I}
(curso de \emph{Matemática}, 365 h) --- convidado a lecionar num curso que não era o seu.


\section*{Proposta de trabalho \textbf{I} --- para engenharia}

\lin{\emph{Headjack} --- interface de imersão sem linguagem intermédia}{proposta aberta}

A única disciplina do mestrado em que não obtive conceito máximo foi \textbf{Interação 3D em
Realidades Mistas} --- e foi a que mais me interessou. O motivo do tropeço é o problema:
\emph{a área não tem formalismo}. Ensina-se por mapas de empatia, entrevistas e método de design,
porque não existe uma linguagem que descreva o que é uma interface entre um corpo e um espaço.
Fica-se com heurísticas onde devia haver uma álgebra.

\medskip
\noindent \textbf{Esta proposta é a resposta a isso}, e vem em duas metades que se sustentam uma à
outra.

\medskip
\lin{1. O formalismo --- e já está escrito}{}
A teoria da dualidade dá o que faltava: um \textbf{par leitura/escrita} com uma involução que os
troca, uma \textbf{régua} por corpo, e uma \textbf{dinâmica} com o seu período. Uma interface deixa
de ser uma coleção de boas práticas e passa a ser um objeto com dois lados e uma passagem entre
eles --- \emph{o que escreve não lê o que escreveu; o que lê não escreve nada que fique}. As
propriedades desejáveis (não dissipar, conservar o invariante, mudar o sítio e não a quantidade)
tornam-se \textbf{verificáveis} em vez de opináveis.

\medskip
\lin{2. O material --- e é produzível}{grafeno + estanho}
O par foi escolhido por medição e não por gosto: o estanho \textbf{fecha o octeto} ($4+4=8$), é
condutor sem ser nobre, e é o \emph{único metal comum com alotropia como a do carbono} --- branco
metálico, cinzento semicondutor. \textbf{Grafeno e estanho são os dois elementos que mudam de
comportamento por rotação estrutural}, e é isso que os torna um par dual e não uma mistura.

\smallskip
\noindent E o compósito \textbf{não interpola: percola}. Abaixo da fração crítica ($\sim0{,}1\%$
para folhas de alta razão de aspeto) o grafeno está disperso e isola; acima, os caminhos ligam-se e
$\sigma \propto (p-p_c)^2$. \textbf{O salto no limiar é de $10^{14}$} --- \emph{o que decide não é a
quantidade, é a geometria de quem toca quem}. \textbf{Isto é o modelo, calculado em}
\texttt{tests/liga.c} --- com a janela de casamento de impedância e a sua dependência da frequência
--- \emph{e não a liga medida em bancada}: essa é precisamente a parte que proponho fazer, e está
dita adiante.

\medskip
\lin{3. A máquina --- e ela existe, corre, e mede-se}{}
Entre as duas metades acima há uma peça construída este mês: uma \textbf{máquina que não guarda
estado e tem uma instrução só}. A lei diz que \emph{toda representação tem dual}, logo é
reversível, logo não precisa de memória. Isso deixou de ser argumento e passou a ser binário:

\begin{center}\small
\begin{tabular}{@{}ll@{}}
\toprule
a instrução & \texttt{MOVE}(destino, sentido) --- e \textbf{mais nenhuma} \\
o estado & \textbf{zero}: sem \emph{runtime}, o binário não tem secção \texttt{.bss} \\
\bottomrule
\end{tabular}
\end{center}

\noindent \textbf{As secções não estão vazias: não existem.} Com \texttt{\_start} próprio e a
saída emitida por \emph{syscall} de dentro, o binário não as tem. E a instrução é uma porque as
outras se \emph{derivam}, não porque se cortaram: os três saltos são um com a \emph{condição} como
argumento; \texttt{LOAD} e \texttt{STORE} são a mesma transferência com o \emph{sentido} trocado;
o salto é a mesma com o \emph{destino} no contador; e a lógica toda deriva do \texttt{NAND},
\emph{que não vive no processador --- vive no slot}.

\smallskip
\noindent \emph{E ela não é uma invenção nossa}: é o mesmo objeto que a tabela de comutação de um
inversor de três níveis, \textbf{descrita atrás no trabalho de conclusão de 2018} --- derivada por
teorema oito anos depois de a ter projetado por outra razão.

\medskip
\lin{4. A inteligência artificial --- onde está, hoje}{em construção, e mede-se}
Esta é a parte \textbf{em aberto}, e digo o ponto exato em vez de a dar por feita. A assistente
existe e corre: \texttt{banco/conversa.c}, \textbf{corpus em disco, sem vocabulário, sem índice
invertido e sem \emph{malloc}} --- a fala cifra-se, desce uma árvore em \texttt{pread}, e a resposta
mora no nó. \textbf{Não chama modelo nenhum de fora}: o sistema é auto-contido, e a bateria não
perde uma unidade sem ele.

\smallskip
\noindent O que se ganhou nesta semana é a \textbf{simetria da procura}, e é a lei aplicada a uma
peça de software. A busca tratava só a fala que tem \emph{a mais} --- havia duas operações
(prefixo e subsequência) e ambas \textbf{tiram} símbolos, nenhuma acrescenta. Era meia dualidade, e
media-se assim: com uma palavra a mais, $252$ de $252$ falas do corpus alcançavam a resposta; com
uma palavra \emph{a menos}, \textbf{zero}. O dual em falta é a \emph{evolução}: a fala cabe inteira,
acaba num nó que não responde, e é o \textbf{banco} que desce o resto sozinho.

\begin{center}\small
\begin{tabular}{@{}lll@{}}
\toprule
 & antes & agora \\
\midrule
uma palavra \emph{a mais} & $252/252$ & $252/252$ \\
uma palavra \emph{a menos} & $0/252$ & $\mathbf{216/252}$ \\
e a resposta é a \textbf{mesma} & --- & $216$ iguais, \textbf{0 diferentes} \\
\bottomrule
\end{tabular}
\end{center}

\noindent \textbf{O número que importa é o zero, não o 216.} Que ela responda mede pouco --- uma
completação que invente também responde. Que responda \emph{exatamente} o que a fala inteira
responderia, $216$ vezes em $216$, é o que mostra que não inventou. E onde o caminho \textbf{bifurca}
ela não escolhe: declara a bifurcação e mostra os ramos, porque a continuação forçada não é um
palpite --- é a única que existe.

\smallskip
\noindent \textbf{O que falta, dito sem rodeios}: o corpus é pequeno e ensinado à mão; a composição
entre falas distintas (fundir dois corpos e não seguir um caminho) está desenhada e \emph{não}
construída; e nada disto foi exposto a uso real fora da bancada. \textbf{É trabalho a fazer, e é
onde procuro colaboração.}

\medskip
\lin{O que procuro para levar isto adiante}{}
\begin{itemize}
\item \textbf{Via software} --- realizar a passagem como transporte de estado entre agentes: o que
um escreve, outro lê, com o invariante verificado em cada travessia. \textbf{Esta parte já existe e
corre}: o \texttt{broca-so} (688 commits, $\sim$355\,000 linhas de código-fonte) orquestra agentes sobre um motor
congelado, e o seu átomo é declarado no cabeçalho --- \emph{«gato $A=[[1,1],[1,0]]$»}, que é a
matriz de Fibonacci, que é a segunda lei da teoria. A sua bateria corre a zero, mas
\textbf{o trabalho não está terminado}: falta a passagem entre agentes ser verificada em produção,
e não só em teste.
\item \textbf{Via física} --- produzir a liga e medir o que o modelo prevê: o limiar de percolação,
o expoente quadrático, e a janela de casamento. \textbf{É a parte que não consigo fazer sozinho},
e é onde procuro laboratório ou parceiro industrial.
\end{itemize}

\medskip
\noindent \emph{A inteligência artificial que fecha o circuito já existe e não é a novidade.} A
novidade é haver \textbf{um formalismo para a interface} e \textbf{um material escolhido por
medição para a realizar} --- e as duas coisas estarem escritas, abertas e verificáveis.

\section*{Proposta de trabalho \textbf{II} --- para jogo e cinema}

\lin{Um mundo construído, e disponibilidade para trabalhar nele com quem o queira levar}{}

O terceiro volume é um romance: \textbf{360 páginas, 26 partes, 146 capítulos, zero fórmulas}. Não
é a teoria explicada --- é a mesma estrutura contada como história, e lê-se sem saber que a outra
existe.

\smallskip
\noindent \textbf{O mundo tem duas metades opostas, e isso foi construído e não acontecido}: a
primeira orgânica --- a terra, a floresta e o mar, o que estava aqui antes de haver reino; a
segunda metálica e espacial --- a Fortaleza, a Alfândega Dimensional, o Telescópio, a Fábrica.
Entre elas, uma armadura que \emph{muda de cor a cada salto}, um cristal que escreve e um traje que
lê. Fecha numa parte chamada \emph{O Zero --- o ponto de simetria}.

\smallskip
\noindent \textbf{E tem uma mecânica, que é o que um romance normalmente não tem.} A estrutura é o
xadrez, a sério e não como metáfora: cada capítulo é dual --- parte branca e parte negra ---
e a semente do mundo é o \textbf{gambito da dama recusado}, com a razão dita em texto corrido:
\emph{quem aceita ganha material e paga com o centro; quem recusa fica com menos e continua largo}.
Há um motor de xadrez próprio no repositório (\texttt{chess}, 819 commits) que joga estas posições.
\textbf{Há aqui um jogo por fazer, e o material de desenho já está escrito.}

\medskip
\lin{O que ofereço, e em que posição}{}
Trago o mundo, a estrutura e a coerência interna --- e a disponibilidade para o trabalhar com quem
souber das partes que eu não sei. \textbf{Aceito direção de arte}, e igualmente \textbf{integrar
equipa artística em qualquer posição}: \emph{direção não é exigência}. Prefiro ver isto feito com
outros a vê-lo intacto e por fazer.

\section*{Proposta de trabalho \textbf{III} --- para registo distribuído}

\lin{\emph{CristalChain} --- um livro-razão cuja \textbf{cifra} conserva}{proposta aberta}

Uma corrente de blocos é, no fundo, uma pergunta sobre \textbf{conservação}: nada se cria do nada,
nada se perde. Hoje isso é garantido por \emph{convenção} --- regras de emissão, consenso entre
participantes, penalizações para quem se desvie. \textbf{A proposta é outra: que a conservação seja
uma propriedade da própria cifra}, e não uma regra imposta por fora a um sistema que, sem ela,
conservaria na mesma.

\smallskip
\noindent E a cifra não foi escolhida --- \textbf{foi obrigada}. Há duas famílias de números
disponíveis para cifrar, e elas fazem coisas opostas: \emph{uma inverte o sinal do que guarda a cada
passo; a outra preserva-o}. Cifrar um livro-razão com a primeira é assentá-lo sobre algo que, por
construção, não conserva.

\smallskip
\noindent Para quem quiser a âncora técnica em vez da imagem: são \textbf{unidades de corpos
quadráticos reais}, e o que as separa é o \textbf{determinante da matriz que gera a recorrência} ---
$-1$ na família conhecida (a de Fibonacci e suas irmãs), $+1$ na outra. A segunda é a que este
trabalho chama \textbf{ouro branco}, e ela distingue-se por não possuir sequer elementos de norma
$-1$ --- isto é, \emph{não tem o lado que anti-conserva}. \textbf{A demonstração está no catálogo},
por dois caminhos independentes que concordam; aqui fica só o motivo, que é o que uma proposta
precisa de dizer.

\medskip
\lin{O contrato não se assina --- liquida-se}{}
Um contrato inteligente tem uma dificuldade conhecida: \textbf{parar}. A solução em uso é comprar a
paragem por fora, dando ao contrato um orçamento de execução --- ele corre até o crédito acabar.
\textbf{Aqui a paragem vem de dentro}: o objeto sobre que o contrato trabalha é finito, e o que é
finito acaba por voltar ao princípio. Não há contador nenhum a vigiá-lo, e mesmo assim ele fecha
--- e fecha \emph{sempre} abaixo de um limite que a própria estrutura anuncia de antemão.

\smallskip
\noindent Daí decorrem as outras duas dificuldades canónicas. \textbf{Quem executa não é escolhido
pela entrada}, mas pela natureza do problema --- portanto uma entrada hostil não consegue escolher
quem corre, que é a raiz das falhas por reentrância. E \textbf{parte do que seria preciso perguntar
ao mundo deriva-se do que já foi fornecido}, sem consultar ninguém --- o que reduz a superfície do
oráculo sem eliminar o problema, e essa distinção fica dita já a seguir.

\smallskip
\noindent A diferença de fundo é simples: \textbf{o contrato assinado} declara cláusulas e alguém
verifica se foram cumpridas; \textbf{o contrato que se liquida} recebe metade e deriva o resto, e
liquidar é apenas correr. Não há contraparte a honrar nem árbitro a julgar --- e a recusa também é
uma liquidação, não uma avaria.

\medskip
\lin{O que existe e o que não existe --- para não haver equívoco}{}
\textbf{O que existe}: a álgebra, os medidores que a verificam e correm a zero, e o papel económico
da CristalChain escrito no corpo da obra --- \emph{``um livro-razão onde nada se forja e nada se
perde''}. \textbf{O que não existe}: rede, nós, consenso entre pares. \textbf{Isto não é uma
blockchain e não resolve dupla despesa}; o que se verifica a cada corrida é teste e não prova
formal; e dados vindos do mundo continuam a precisar de uma fonte de confiança.

\smallskip
\noindent \emph{O que trago é a razão pela qual a cifra é aquela e não outra} --- e procuro quem
saiba da parte distribuída, que não é a minha. \textbf{A demonstração está no catálogo}
(\emph{A CristalChain --- um livro-razão cuja cifra conserva}) e os medidores no repositório, ambos com endereço no fim deste
documento.

\section*{Como trabalho}

Digo isto porque um parceiro precisa de saber, e porque está escrito nos meus registos de qualquer
forma.

\medskip
\noindent \textbf{Aprendo sozinho e trabalho por imersão.} Na graduação estudava em casa e ia às
aulas o indispensável --- o histórico mostra as duas coisas: notas máximas nas disciplinas teóricas
e reprovações por frequência nas de laboratório. Não é irregularidade: é que a minha presença é
função do conteúdo. Nos dois cursos de análise que escolhi fazer, a frequência foi de 100\,\%.

\medskip
\noindent \textbf{Meço em vez de argumentar.} Toda afirmação minha vem com um teste que a pode
derrubar, e uma parte considerável do meu trabalho consiste em derrubar as minhas próprias
conclusões. Prefiro um resultado negativo publicado a um positivo por verificar.

\medskip
\noindent \textbf{O que me falta, e o que procuro.} Trabalho melhor com problemas difíceis e prazo
longo do que com rotina; e o trabalho ganharia com quem discuta o formalismo de perto --- em
álgebra, em topologia, ou em quem o queira aplicar. \textbf{É isso que procuro.}

\section*{A colaboração que procuro tem dois lados}

E digo os dois, porque \textbf{são um par e não uma alternativa com uma opção a sério e outra por
educação}. Um precisa de mim presente; o outro funciona sem mim --- e é isso que os torna duais e
não graus da mesma coisa.

\medskip
\noindent \textbf{Presencial --- com equipa e estrutura.} É onde estão as partes que não consigo
fazer sozinho: produzir a liga e medir o que o modelo prevê, levar a interface a um laboratório,
pôr a passagem entre agentes a correr em produção e não só em teste. Aqui procuro
\textbf{grupo, bancada e prazo}, e aceito qualquer posição em que seja útil.
\textbf{Se é isto que lhe interessa, escreva-me} --- o contacto está no topo, e respondo a todas.

\medskip
\noindent \textbf{Remota --- e sem me pedir licença.} Está tudo publicado e aberto: quem se
interessar \textbf{pega no conteúdo e avança}, sem combinar nada comigo antes. Não é preciso
autorização, não é preciso avisar, e \emph{não é preciso devolver} --- a licença já o permite e a
intenção é essa. \textbf{Se depois quiser falar comigo, falo com muito gosto; se não quiser, o
trabalho serviu na mesma.} O que não quero é que fique parado à espera de mim.

\medskip
\lin{Onde está tudo}{}
\begin{center}\footnotesize
\begin{tabular}{@{}llr@{}}
\toprule
documento & endereço & extensão \\
\midrule
a teoria --- o que explica & \texttt{goldenkingdom.patriatechnology.com/docs/teoria.pdf} & 58\,pp. \\
o catálogo --- o que verifica & \texttt{goldenkingdom.patriatechnology.com/docs/catalogo.pdf} & 446\,pp. \\
o enredo --- o que conta & \texttt{goldenkingdom.patriatechnology.com/docs/enredo.pdf} & 360\,pp. \\
os três num volume & \texttt{goldenkingdom.patriatechnology.com/docs/livro.pdf} & 867\,pp. \\
\midrule
o código, clonável & \texttt{git clone https://goldenkingdom.patriatechnology.com/repo.git} & \\
\bottomrule
\end{tabular}
\end{center}

\noindent \textbf{O repositório é o mesmo que gera os documentos}, e traz a bateria inteira: as
afirmações dos três volumes têm um medidor que as pode derrubar, e correm todos com um comando.
\emph{Quem quiser verificar antes de acreditar --- que é o que eu faria --- tem com quê.}

\vfill
{\footnotesize\color{cinza}\noindent Este documento e todo o trabalho aqui descrito estão sob
\textbf{Creative Commons Atribuição 4.0} (textos) e \textbf{MIT} (código). Use, adapte, publique
--- basta citar a origem. \\
\texttt{goldenkingdom.patriatechnology.com}\par}

\end{document}
